\section{引言}

\subsection{问题陈述}

在有限信息与有限分辨率约束下，“观测”不能被等同为对连续状态的全量访问，而应被视为一个由几何参数、动力学采样与读出核共同决定的协议。若要求该协议满足：（i）本体层允许以高维周期对象实现非周期结构；（ii）观察者差异可被严格参数化；（iii）读出过程产生可验证的不变量；（iv）稳定态由可计算的规范形给出；（v）时间的出现可被解释为信息精化而非外加参数，则需要一条从几何—动力学—信息论到可观测量的闭合推理链。

\subsection{主要贡献}

本文在严格区分\textbf{定义}、\textbf{可证明定理}、\textbf{模型类假设}与\textbf{解释层命题}的规则下，完成如下工作：
\begin{enumerate}
  \item \textbf{（定义层）}以六维晶格切投影方案构造三维准周期点集，并以内部相位参数化观察者，实现“同一本体—多实现”的可审计表述。
  \item \textbf{（定理层）}以环面旋转一维因子与二分割读出得到 Sturmian 机械词，并给出复杂度最小性不变量。
  \item \textbf{（定义层）}构造 Zeckendorf 合法语言与稳定化核（folding），将任意二值读出窗口规范化至“无相邻 $1$”的稳定态空间。
  \item \textbf{（定理层 + 模型假设）}定义时间纤维为稳定化核的预像，证明在遍历与生成分割条件下观测历史柱集测度的指数收缩率由度量熵给出，从而将“时间箭头”锚定在可计算的信息率上。
  \item \textbf{（模型定义）}以节点—纤维入射图刻画观测可达闭包，并给出由纤维统计诱导的“类质量/类延迟”等可观测量函数子。
  \item \textbf{（解释层命题）}提出可见性投影与黄金比标度的固定点框架，讨论非整数有效维与算术接口，明确其不进入闭层定理体系。
\end{enumerate}

\subsection{文章结构}

第 2 章给出切投影方案与观察者参数化；第 3--4 章建立有限资源下的采样与 Sturmian 二值读出；第 5 章定义 Zeckendorf 稳定化核；第 6 章将时间形式化为预像纤维的历史细化并给出信息率刻画；第 7 章给出 Levi 入射图闭包与可观测量函数子；第 8 章讨论遍历下的统计等价；第 9 章给出多尺度与有效维讨论；第 10 章给出可复现协议；附录给出证明要点、算法与扩展接口。

